\section{Introduction}
\label{sec:intro}

We develop \emph{Thurstone portfolios}, a long-only allocation scheme whose
weights are the winning probabilities of a race among assets: each asset draws a
noisy performance, and its weight is the probability that it comes out best. The
race has a second reading as a Thurstonian model of choice, which we set against
the passive industry's capitalization weighting and its implicit appeal to Luce's
choice axiom. That Lucian habit is only superficially supported by the Capital
Asset Pricing Model (CAPM): what the CAPM endorses is the efficiency of the
whole-universe market portfolio, a non-Lucian object no investor on a restricted
universe actually holds (\Cref{sec:capm}).

The reason to read allocation as a Thurstonian rather than a Lucian choice is a
single property: \emph{redundancy consistency}. Add a second asset that behaves
almost exactly like one already held, and a sensible rule should \emph{split} the
position across the two near-duplicates and leave the rest of the book untouched;
it must not let the duplicated exposure count twice. This is the portfolio form of
the red-bus/blue-bus problem that unseated the logit model in discrete choice
\parencite{mcfadden1974}. A Lucian rule---capitalization weighting included---fixes
the ratio of any two weights independently of the rest of the field, so inserting a
near-duplicate inflates the duplicated exposure and drains weight from everything
else (\Cref{prop:restriction,prop:redundancy}). A Thurstonian race does not: two
assets that move together compete with \emph{each other}, and in the limit a
co-moving pair collapses to a single horse that holds their combined weight.

Covariance-based allocators are not blind to this; they reach the same
de-duplication, but lurchily. Mean--variance settles the split between near-
duplicates by inverting $\Sigma$, and the direction in which two assets become
identical is exactly the direction in which $\Sigma$ loses rank; so the population
answer is right while its sample estimate, dividing through a vanishing eigenvalue,
swings between the twins and flips sign with the data \parencite{markowitz1952}.
The inversion-free clustering allocators (hierarchical risk parity, the
Schur-complementary family) de-duplicate by grouping instead, but discretely: the
split jumps as cluster membership flips. The race de-duplicates \emph{smoothly},
with a closed characterization (\Cref{prop:redundancy}) and without ever forming
$\Sigma^{-1}$.

That is the organizing idea of the paper. Redundancy is near-singularity, so three
properties usually treated apart turn out to be one. The race is robust to
ill-conditioning because it samples from a correlation rather than inverting one
(\Cref{sec:robust}); it turns over little because its weights are locally Lipschitz
in the correlation and move only as far as the correlation moves, where the
clustering methods re-seriate each period (\Cref{thm:smooth}, \Cref{sec:experiments});
and it splits redundant exposure gracefully. The three are one behaviour---moving
smoothly along the direction in which assets become collinear and a covariance
inverse would explode.

One property genuinely stands apart, because it is not about conditioning at all.
Unlike any covariance functional, the race admits \emph{tail dependence}: driven by
a tail-dependent simulation it responds to lower-tail co-movement that no second
moment can encode, so a cluster that crashes together---but only in the
crash---de-duplicates in the tail (\Cref{prop:tail}). The winner of a race is an
order statistic of the joint draw, and an order statistic reads the tail copula
where a covariance cannot.

Read the other way, a winning probability is a credit as much as a weight: a
redundancy-aware attribution of credit among correlated contributors. The same
machinery combines forecasts, where stable combinations are known to beat the
least-squares-optimal one, and we report a first study (\Cref{sec:combination}).

These rest on properties established below: feasibility on the simplex with
neither optimization nor inversion, exact reproduction of the benchmark when the
reference and target laws agree, redundancy consistency and diversification
monotonicity, tail consistency, an implied mean--(convex-regularizer) objective
$w = \nabla G_S(\theta)$ for any performance law $S$, and a Lipschitz bound on
turnover (\Cref{sec:properties,sec:smoothness}).

\section{Background}
\label{sec:background}

\subsection{Why the CAPM does not rescue naive indexing}
\label{sec:capm}

The passive industry defends capitalization weighting by appeal to the CAPM. We
grant the CAPM in full; the point is that its support for the practice is
superficial. What the model endorses is the efficiency of the whole-universe
market portfolio, a non-Lucian object the indexer never holds, and that
endorsement does not transfer to the restricted universes where weighting
decisions are made.

\paragraph{What the CAPM says.}
Under the CAPM every investor is a mean--variance optimizer who agrees on the
expected returns $\mu$ and on the \emph{same} covariance $\Sigma$. Two-fund
separation then has each investor hold the risk-free
asset together with the \emph{same} tangency portfolio, proportional to
$\Sigma^{-1}(\mu - r\mathbf{1})$. Since all investors hold the identical risky
portfolio, market clearing forces it to equal the supply, the
capitalization-weighted market, so $\mathrm{cap} \propto \Sigma^{-1}(\mu -
r\mathbf{1})$ and the market portfolio is efficient.

Capitalization weighting is optimal precisely \emph{because the covariance is
shared and the universe is the
whole world}: it is the one portfolio on which all investors, reasoning from a
single $\Sigma$ over all assets, agree.\footnote{One wrinkle we do not dwell on:
investors have a well-established motivation to shrink the covariance
\parencite{ledoitwolf2004}, and an aggregate of heterogeneous shrunk estimates
need not clear to a mean--variance-efficient market portfolio, so even the
equilibrium claim is not airtight.}

\paragraph{Choice and efficiency.}
Given these imperfections it is natural to ask whether allocation admits other,
less normative models. The choice literature offers a rich menu, and in particular
two accounts long viewed as rivals in the heyday of psychometric choice theory:
Thurstone's law of comparative judgment \parencite{thurstone1927} and the Lucian
(logit) model. They part on whether the relative standing of two options may
depend on the rest of the field, and that fault line is the one that matters here.

Write $w_i(\mathcal U)$ for the weight a rule assigns to asset $i$ within a
universe $\mathcal U$. Capitalization weighting sets $w_i(\mathcal U) =
\mathrm{cap}_i / \sum_{j\in\mathcal U}\mathrm{cap}_j$, so the ratio
$w_i(\mathcal U)/w_j(\mathcal U) = \mathrm{cap}_i/\mathrm{cap}_j$ does not depend on
$\mathcal U$. That is Luce's choice axiom \parencite{luce1959}, independence of
irrelevant alternatives, and we call such a rule \emph{Lucian}. Efficiency has no
such invariance: the tangency weights on $\mathcal U$ are $w^{*}(\mathcal U)
\propto \Sigma_{\mathcal U}^{-1}(\mu_{\mathcal U} - r\mathbf{1})$, whose ratios turn
on the entire inverse covariance $\Sigma_{\mathcal U}^{-1}$ and move when an asset
enters or leaves $\mathcal U$ (unless $\Sigma_{\mathcal U}$ is diagonal). The CAPM
makes the two objects coincide at a single universe, the whole world, where
clearing fixes $\mathrm{cap}\propto\Sigma^{-1}(\mu-r\mathbf 1)$.

\begin{proposition}[Lucian weighting is inconsistent under restriction]
\label{prop:restriction}
Capitalization weighting is the Lucian rule with values $\mathrm{cap}_i$. If
returns are correlated (non-diagonal $\Sigma$), then generically a Lucian rule
coincides with mean--variance efficiency on at most one universe in a nested
family. Under the CAPM that universe is the full asset world $U$; on any strict
sub-universe $S \subsetneq U$ the capitalization-weighted portfolio is generically
inefficient, and the discrepancy is governed by the heterogeneity of the assets'
covariances with the excluded set $U \setminus S$.
\end{proposition}
\begin{proof}
Partition the market $U = S \cup X$ into the retained assets $S$ and the excluded
$X = U\setminus S$, with market weights $m = (m_S, m_X)$. CAPM equilibrium reads
$\mu - r\mathbf 1 = \gamma\,\Sigma_U\, m$ for some $\gamma > 0$; its $S$-block is
$\mu_S - r\mathbf 1 = \gamma(\Sigma_{SS} m_S + \Sigma_{SX} m_X)$, so the
restricted-universe tangency weights are
\begin{equation*}
  q_S \;\propto\; \Sigma_{SS}^{-1}(\mu_S - r\mathbf 1)
  \;=\; \gamma\big(m_S + \Sigma_{SS}^{-1}\Sigma_{SX}\, m_X\big),
\end{equation*}
while restricted capitalization weighting uses $m_S$. The two coincide iff
$\Sigma_{SS}^{-1}\Sigma_{SX}\, m_X \propto m_S$, which holds generically only at the
full universe ($X = \varnothing$), fixing the single coincidence point. The wedge
$\Sigma_{SS}^{-1}\Sigma_{SX}\, m_X$---the regression of the retained assets on the
excluded ones---is precisely the part of each cap weight that reflects covariance
with $U\setminus S$, irrelevant to the $S$-frontier yet baked into the full-market
weight.
\end{proof}

The CAPM's support for the Lucian habit is therefore the single-point coincidence
of \Cref{prop:restriction}: it endorses cap weighting on the one universe nobody
trades, and is silent, indeed contrary, everywhere else. This is the
choice-theoretic form of Roll's critique \parencite{roll1977}: the only testable
content of the CAPM is the efficiency of the \emph{true} market portfolio, and no
proxy on a restricted universe inherits it \parencite{prono2009}. We do not
contradict the CAPM, only decline to extend it past its premises. Fundamental
indexing reaches cap weighting's sub-optimality by a different route, price
inefficiency \parencite{hsu2006}, while shrinkage \parencite{ledoitwolf2004} and
tracking-error analysis \parencite{roll1992} document the fragility of the inputs.

\subsection{Asset allocation as a choice}
\label{sec:choice}

If capitalization weighting is a Lucian \emph{habit} rather than a principle, what
is an \emph{informationally modest} alternative to naive indexing? We read
allocation as a \emph{choice}. The market, in aggregate, selects among assets, and
a portfolio weight is the probability with which an asset is chosen. Thurstone's
law of comparative judgment \parencite{thurstone1927} is the canonical model of
such a choice: each alternative $i$ carries a latent quality, an \emph{ability}
$a_i$, and a noisy \emph{performance} $X_i = a_i + \varepsilon_i$, and the
alternative with the extreme performance is chosen. The probability that $i$ is
chosen depends on the whole field and on the dependence among the performances,
the property the Lucian (logit) model forbids by axiom. Two assets that move
together compete with each other, not independently, and a choice model charges for
that where a Lucian one cannot.

Using these choice probabilities as weights is the construction of this paper:
\Cref{sec:method} fixes the field's abilities by \emph{calibration} to a chosen
benchmark, then \emph{tilts} by re-running the choice under the dependence actually
realised.

The tilt is informationally modest.
It needs only that dependence, no return forecast, and it \emph{anchors} to the
benchmark, departing from the prevailing allocation only as far as the correlation
warrants. Treating the market as at least weakly efficient is the prudent default;
discarding that anchor to bet against the index on thin evidence is the error that
undid many entrants in the M6 forecasting competition.

\subsection{Thurstone's eclipse, and the Fast Ability Transform}
\label{sec:fat}

The Thurstonian (probit) choice model is older and strictly more general than the
Lucian (logit) one, and it does not assume independence of irrelevant
alternatives. It lost the applied field on a single axis, tractability. With more
than two alternatives the choice probabilities are Gaussian orthant integrals with
no closed form, and the \emph{inverse} problem, recovering latent abilities from
observed choice shares, was considered prohibitive. Logit has closed-form
probabilities and a trivial inverse, and so came to dominate discrete choice,
ranking, and learning-to-rank, carrying its IIA assumption with it.

Two developments lift the obstruction here. First, the forward race is cheap to
\emph{simulate}: the winning frequencies of $10^4$ to $10^5$ correlated draws cost
milliseconds, and a common-seed construction (\Cref{sec:smoothness}) makes the
simulation \emph{smooth} in the dependence, so it can be re-used online without
churning the weights. Second, and decisively for calibration, the \emph{Fast
Ability Transform} of \textcite{cotton2021} evaluates the forward and inverse maps
of the independent race in \emph{linear time} on a lattice, convolving competitor
densities and tracking tie multiplicity exactly, so the once-prohibitive inverse
from target weights to abilities is now $O(n)$ and Monte-Carlo-free. A theoretical
curiosity becomes a practical allocator.

The transform is the hinge of everything that follows. We develop several uses of
the apparatus, with different benchmarks to anchor to, different reference
correlations, and target simulations carrying whatever dependence one wishes to
match, but each routes through this one calibration map. None is practical without
a fast, exact inverse from weights to abilities; with it, the rest is sampling.

\subsection{Avoiding matrix inversion}
\label{sec:robust}

The operational fragility of mean--variance optimization is the inversion of an
estimated covariance \parencite{markowitz1952}: as $\Sigma$ nears singularity,
$\Sigma^{-1}$ and with it the weights explode. The discipline's responses are
shrinkage toward a structured target \parencite{ledoitwolf2004} and, more
radically, allocators that never form $\Sigma^{-1}$ at all (hierarchical risk
parity, the Schur-complementary family, risk parity), trading point optimality for
stability and scale. The Thurstone race belongs to this inversion-free family: it
\emph{samples} from a correlation rather than inverting one, so it stays well posed
even when $\Sigma$ is singular or rank-deficient.

Robustness to \emph{conditioning} is not robustness to \emph{tails}. Every method
just named is a functional of the covariance, or of a correlation derived from it.
A covariance is a second moment; it summarises how assets move together on average
and says nothing about how they move together in the extreme. Two universes with
identical covariance, hence identical minimum-variance, HRP, Schur, and
risk-parity weights, can carry entirely different probabilities of a simultaneous
crash. The quantity that separates them, tail dependence, is not a moment, and no
re-estimation or shrinkage of $\Sigma$ recovers it.

\subsection{The tail a covariance cannot see}
\label{sec:tailmotiv}

Concretely: how often do the Dow's constituents fall \emph{together}, and by how
much? A multivariate Gaussian fit to their daily returns, pinning the means and the
full covariance, reproduces a shallow joint decline well enough, but its joint tail
thins far too fast. All of them closing down on the same day is only mildly
underpredicted; deepen the event and the gap widens without bound. Over
2010--2024, days on which every name fell by more than two percent occurred about
eighty times more often than the fitted Gaussian allows, and beyond three percent
the Gaussian effectively never produces such a day, yet two did. The direction of
the asymmetry is the fingerprint: the same fit underpredicts joint \emph{crashes}
far more than joint \emph{rallies} (\Cref{tab:tail}).\footnote{Daily log returns of
the 28 Dow names with full 2010--2024 history (Yahoo Finance; average pairwise
correlation $0.42$), against a Gaussian fit to their mean and covariance, the
probabilities by $10^{8}$ Monte-Carlo draws. The deepest rows rest on two to four
days, so the precise multiples are noisy; the monotone, asymmetric divergence is
not.} The Gaussian gets the average right and the joint extreme catastrophically
wrong, because the lower-tail dependence
\begin{equation}
  \lambda_L \;=\; \lim_{u\downarrow 0}\Prob\!\big(U_i < u \,\big|\, U_j < u\big)
  \label{eq:lambdaL}
\end{equation}
is zero for the Gaussian copula and substantial in markets. No shrinkage or
re-estimation of $\Sigma$ moves $\lambda_L$ off zero; it is invisible to second
moments.

\begin{table}[h]
  \centering
  \caption{All $28$ Dow constituents moving by more than $q$ on the same day
  (2010--2024): empirical frequency, the frequency under a Gaussian fit to the
  mean and covariance, and their ratio. The Gaussian undercounts joint
  \emph{crashes} ever more steeply with depth, and far more than joint
  \emph{rallies}---the lower-tail dependence a second moment cannot carry.}
  \label{tab:tail}
  \begin{tabular}{lccc@{\qquad}ccc}
    \toprule
    & \multicolumn{3}{c}{all down by ${>}q$} & \multicolumn{3}{c}{all up by ${>}q$} \\
    \cmidrule(lr){2-4}\cmidrule(lr){5-7}
    $q$ & emp. & Gaussian & ratio & emp. & Gaussian & ratio \\
    \midrule
    $1.0\%$ & $0.32\%$ & $0.080\%$  & $4\times$         & $0.24\%$ & $0.111\%$  & $2\times$  \\
    $1.5\%$ & $0.21\%$ & $0.012\%$  & $18\times$        & $0.13\%$ & $0.018\%$  & $7\times$  \\
    $2.0\%$ & $0.11\%$ & $0.0014\%$ & $77\times$        & $0.05\%$ & $0.0022\%$ & $25\times$ \\
    $3.0\%$ & $0.05\%$ & ${\sim}0$  & ${>}10^{4}\times$ & $0\%$    & ${\sim}0$  & ---        \\
    \bottomrule
  \end{tabular}
\end{table}

A choice among assets reads the joint extreme directly. The winner of a race is an
\emph{order statistic} of the joint draw, and the law of an order statistic is set
by the tail copula, not by the covariance. Drive the Thurstonian race with a
simulation that carries lower-tail dependence and the weights respond to
$\lambda_L$: assets that crash together are recognised as the near-duplicates they
become in the tail, and their shares collapse to one, which no covariance-based
allocator can arrange (\Cref{prop:tail}). With this in view we make the
construction precise.
